\chapter{Datensicherheit -- Bedrohungen und Schutzmaßnahmen}
\renewcommand{\chaptertitle}{Datensicherheit -- Bedrohungen und Schutzmaßnahmen}

\lehead[]{\normalfont\sffamily\hspace*{-2.00cm}\textcolor{white}{\colorbox{lightblue}{\makebox[1.60cm][r]{\thechapter}}}\hspace{0.17cm}\textcolor{lightblue}{\chaptertitle}}
\rohead[]{\textcolor{lightblue}{\chaptertitle}\normalfont\sffamily\hspace*{0.17cm}\textcolor{white}{\colorbox{lightblue}{\makebox[1.60cm][l]{\thechapter}}}\hspace{-2.00cm}}
%\chead[]{}
\rehead[]{\textcolor{lightblue}{AvHG, Inf, My}}
\lohead[]{\textcolor{lightblue}{AvHG, Inf, My}}

\section{Datensicherheit, Datenschutz und Schutzziele}

Im Alltag werden die Begriffe \emph{Datenschutz} und \emph{Datensicherheit} oft
verwechselt, beschreiben aber zwei unterschiedliche Perspektiven:

\begin{compactitem}
\item \textbf{Datenschutz} schützt das Recht jedes Einzelnen, selbst darüber
zu entscheiden, wer welche personenbezogenen Daten über ihn erheben und
verarbeiten darf (informationelle Selbstbestimmung, z.\,B.\ geregelt in der
DSGVO).
\item \textbf{Datensicherheit} (Informationssicherheit) umfasst alle
technischen und organisatorischen Maßnahmen, um Daten und IT-Systeme --
unabhängig davon, ob sie personenbezogen sind -- vor Verlust, Diebstahl oder
Manipulation zu schützen.
\end{compactitem}

Jede Maßnahme der Datensicherheit lässt sich einem der drei klassischen
\textbf{Schutzziele} (auch \emph{CIA-Trias} genannt) zuordnen:

\begin{description}
\item[Vertraulichkeit (\emph{Confidentiality}):] Daten dürfen nur von
autorisierten Personen eingesehen werden (Schutz vor unbefugtem Mitlesen oder
Datendiebstahl).
\item[Integrität (\emph{Integrity}):] Daten müssen vollständig und unverfälscht
bleiben; unbemerkte Änderungen oder Manipulationen müssen verhindert bzw.\
erkannt werden.
\item[Verfügbarkeit (\emph{Availability}):] IT-Systeme und Daten müssen für
berechtigte Nutzer jederzeit zuverlässig erreichbar und funktionsfähig sein.
\end{description}


\section{Typische Angriffsvektoren im Überblick}

Die Methoden von Angreifern entwickeln sich ständig weiter. Um dennoch den
Überblick zu behalten, lassen sich moderne Angriffe nach ihrem Ansatzpunkt in
vier grundlegende Kategorien einordnen.

\subsection{Der Faktor Mensch: \emph{Social Engineering}}

Viele erfolgreiche Angriffe beginnen nicht mit einem technischen Hack, sondern
mit der gezielten Täuschung von Menschen:

\begin{compactitem}
\item \textbf{Phishing \& Spear-Phishing:} Über täuschend echt gestaltete
E-Mails, SMS (\emph{Smishing}) oder QR-Codes (\emph{Quishing}) werden Opfer auf
gefälschte Anmeldeseiten gelockt, um Passwörter oder Bankdaten preiszugeben.
Beim \emph{Spear-Phishing} sind die Nachrichten gezielt auf eine bestimmte
Person oder Firma zugeschnitten.
\item \textbf{CEO-Fraud \& Deepfakes:} Angreifer geben sich per E-Mail oder
mithilfe KI-generierter Stimm- und Videoimitationen (\emph{Deepfakes}) als
Vorgesetzte oder Bekannte aus, um eilige Überweisungen oder die Herausgabe
vertraulicher Daten zu erwirken.
\end{compactitem}

\subsection{Schadsoftware (\emph{Malware})}

Unter \emph{Malware} (engl.\ \emph{malicious software}) versteht man Programme,
die entwickelt wurden, um Schaden anzurichten oder unbefugten Zugriff zu
erlangen:

\begin{compactitem}
\item \textbf{Ransomware (Verschlüsselungstrojaner):} Eine der größten
aktuellen Bedrohungen für Unternehmen, Behörden, Krankenhäuser und Schulen.
Die Schadsoftware verschlüsselt sämtliche erreichbaren Dateien und
Netzwerklaufwerke und fordert ein Lösegeld (\emph{Ransom}) für den
Entschlüsselungsschlüssel -- oft verbunden mit der Drohung, zuvor gestohlene
Daten zu veröffentlichen (\emph{Double Extortion}).
\item \textbf{Trojaner \& Infostealer:} Programme, die sich als nützliche
Software tarnen, im Hintergrund aber heimlich Zugangsdaten, gespeicherte
Browser-Passwörter oder Sitzungs-Cookies auslesen und an Angreifer senden.
\item \textbf{Computerwürmer:} Schadprogramme, die sich -- anders als klassische
Viren, die eine Wirtsdatei benötigen -- selbstständig über Sicherheitslücken im
Netzwerk von Rechner zu Rechner verbreiten.
\end{compactitem}

\subsection{Angriffe auf Netzwerke und Benutzerkonten}

\begin{compactitem}
\item \textbf{DDoS-Angriffe (\emph{Distributed Denial of Service}):} Tausende
ferngesteuerte, mit Schadsoftware infizierte Rechner (ein sogenanntes
\emph{Botnetz}) richten gleichzeitig massenhaft Anfragen an einen Server oder
eine Website. Der Dienst wird überlastet und bricht zusammen -- eine gezielte
Verletzung des Schutzziels \emph{Verfügbarkeit}.
\item \textbf{Man-in-the-Middle-Angriff (MitM):} Ein Angreifer schaltet sich
(z.\,B.\ in einem ungesicherten öffentlichen WLAN) unbemerkt in die
Kommunikation zwischen Client und Server ein, um unverschlüsselte Daten
mitzulesen oder zu verändern.
\item \textbf{Brute-Force \& Credential Stuffing:} Beim \emph{Brute-Force}-Angriff
werden Passwörter systematisch durchprobiert. Beim noch häufigeren
\emph{Credential Stuffing} nutzen Angreifer Zugangsdaten aus früheren
Datenlecks anderer Websites automatisiert aus, weil viele Menschen dasselbe
Passwort mehrfach verwenden.
\end{compactitem}

\subsection{Angriffe auf Software und Webanwendungen}

\begin{compactitem}
\item \textbf{SQL-Injection \& Cross-Site-Scripting (XSS):} Werden
Benutzereingaben in Webanwendungen oder Datenbankprogrammen nicht sauber
geprüft oder maskiert, können Angreifer eigenen SQL-Code (siehe Kapitel~37)
oder schädlichen JavaScript-Code einschleusen.
\item \textbf{Zero-Day-Exploits:} Das gezielte Ausnutzen einer
Sicherheitslücke in einem Betriebssystem oder Programm, die dem Hersteller noch
unbekannt ist (es gibt also noch $0$ Tage lang ein Sicherheitsupdate).
\item \textbf{Supply-Chain-Angriffe:} Statt das eigentliche Ziel direkt
anzugreifen, manipulieren Angreifer einen Zulieferer -- etwa eine verbreitete
Programmbibliothek oder den Update-Server einer Softwarefirma --, sodass das
Schadprogramm über ein scheinbar offizielles Update verteilt wird.
\end{compactitem}


\section{Grundlegende Schutzmaßnahmen}

Gegen die genannten Bedrohungen hilft keine einzelne Maßnahme, sondern nur ein
mehrstufiges Sicherheitskonzept:

\begin{compactitem}
\item \textbf{Starke Authentifizierung:} Für jeden Dienst ein eigenes, langes
Passwort (verwaltet über einen \emph{Passwort-Manager}) sowie konsequente
Nutzung der \emph{Mehr-Faktor-Authentifizierung (MFA)}, sodass ein gestohlenes
Passwort allein nicht ausreicht.
\item \textbf{Zeitnahe Updates (Patch-Management):} Betriebssysteme, Browser
und Anwendungen regelmäßig aktualisieren, um bekannte Sicherheitslücken zu
schließen.
\item \textbf{Regelmäßige Backups (3-2-1-Regel):} Der wirksamste Schutz gegen
Datenverlust und Ransomware: mindestens \textbf{3} Kopien der Daten auf
\textbf{2} unterschiedlichen Speichermedien, davon mindestens \textbf{1} Kopie
räumlich getrennt und offline (vom Netzwerk getrennt, damit Ransomware das
Backup nicht mitverschlüsseln kann).
\item \textbf{Verschlüsselung:} Konsequente Nutzung verschlüsselter
Verbindungen (z.\,B.\ HTTPS/TLS) sowie Verschlüsselung sensibler Datenträger.
\item \textbf{Prinzip der minimalen Rechte (\emph{Least Privilege}):} Jeder
Benutzer und jedes Programm erhält nur genau die Zugriffsrechte, die für die
jeweilige Aufgabe unbedingt erforderlich sind (im Alltag nicht als
Administrator arbeiten!).
\end{compactitem}


\section{Aufgaben und Diskussionsfragen}

\begin{compactenum}[a)]
\item Ordne die folgenden Vorfälle jeweils den verletzten Schutzzielen
(\emph{Vertraulichkeit}, \emph{Integrität}, \emph{Verfügbarkeit}) zu und
begründe deine Entscheidung:
\begin{compactenum}[1.]
  \item Ein Online-Shop ist wegen eines DDoS-Angriffs sechs Stunden lang nicht
  erreichbar.
  \item Ein Angreifer ändert mithilfe einer SQL-Injection den Kontostand in
  einer Bankdatenbank.
  \item Eine Liste mit Kundenpasswörtern wird gestohlen und im Internet
  veröffentlicht.
  \item Eine Schulverwaltung wird Opfer eines Ransomware-Angriffs und kann nicht
  mehr auf ihre Dateien zugreifen.
\end{compactenum}

\item Erkläre, warum ein Backup auf einer dauerhaft am Computer angeschlossenen
externen Festplatte gegen moderne Ransomware oft nicht ausreicht.

\item Überlege, warum das Wiederverwenden desselben Passworts bei mehreren
Online-Diensten auch dann gefährlich ist, wenn das eigene Passwort sehr lang
und komplex ist.
\end{compactenum}
